\documentclass[12pt]{article}
\usepackage[utf8]{inputenc}
\usepackage[T1]{fontenc}
\usepackage{hyperref}
\title{Zumaya Language Resources for K-12 Curriculum Development}
\author{Compiled by Claude Sonnet 5.5 (claude-sonnet-5-5) for Andr\'as Kornai}
\date{}
\begin{document}
\maketitle

\section{Introduction}
Zumaya (ISO 639-3 zuy; Glottocode zuma1239) is an extinct Chadic (Afro-Asiatic) language of northern Cameroon, in the Masa group. It is known from only about ten words recorded by Daniel Barreteau from what was probably the last speaker. Speakers are reported to have shifted to Fulfulde. Sources disagree on when it died out. This is far too little to build a curriculum on.

\section{Primary Educational Resources}
None exist as far as the searches showed.

\section{Linguistic and Reference Materials}
\begin{itemize}
\item Wikipedia entry on the Zumaya language. \url{https://en.wikipedia.org/wiki/Zumaya_language}
\item Glottolog entry for Zumaya, with its classification and status. \url{https://glottolog.org/resource/languoid/id/zuma1239}
\item Catalogue of Endangered Languages entry for Zumaya. \url{https://endangeredlanguages.com/lang/559}
\end{itemize}

\section{Religious and Community Resources}
None were found.

\section{Digital and Audio Resources}
None were found.

\section{Curriculum Development Context}
There is no known corpus, orthography or community of speakers. A heritage programme could at most teach about the language and its history, and the neighbouring Masa-group languages and local Fulfulde could be used for comparison. A real revitalization curriculum is not feasible with the documentation found.

\end{document}